% Consolidated guide-selftest09 from 283f4ad03ab2102bd6bfdf75a06052da75ec7dbf; source fragments preserved.
% 2026-11-24 八批严格二校第8批：历史片段哈希对应修订前版本。

% === BEGIN guide-selftest09-choices.tex blob 66928b4f4a1ed5aad3847c87d9b75df9d08c30b9 ===
% 《大学物理学学习指导》第5版，自测题9（光的干涉）选择题9.1.1—9.1.11
% 原书印刷202—204页，PDF物理210—212页；原图依文字方向重述。
\subsection{自测题9：干涉选择题}

\sourcetest{9.1.1}{202}{210}
\problemtext 单色光真空波长为$\lambda$，在折射率$n$的媒质中
从A传播到B，两点光相位差为$3\pi$，求该段光程。
\workings
假定相位差为沿传播方向累积的未折叠相位延迟（并非仅知模$2\pi$的相位），光程与相位差的关系为
$\Delta\phi=2\pi\mathcal L/\lambda$，故
\[
\boxed{\mathcal L=\frac{\Delta\phi}{2\pi}\lambda=\frac32\lambda}.
\]
原书选\textbf{A}。
\crosscheck 若只知相位差模$2\pi$为$\pi$，还允许$\mathcal L=(j+1/2)\lambda$（$j=0,1,\ldots$）。问的是光程$n\,AB$，不是几何长度$AB=3\lambda/(2n)$。

\sourcetest{9.1.2}{202}{210}
\problemtext 图9.1.1：透射薄膜折射率$n_2$、厚$e$，上下环境
$n_1<n_2>n_3$；入射光在介质$n_1$中的波长为$\lambda_1$。
两束反射光的光程差是多少？
\workings
在正入射、透明非磁介质的条件下，上表面反射从疏到密，出现一次$\pi$相变；
下表面从密到疏无半波相变。
真空波长$\lambda_0=n_1\lambda_1$，反射两路的
有向有效光程差可记为
\[
\boxed{\Delta=2n_2e-\frac{n_1\lambda_1}{2}},
\]
对应原书选\textbf{C}。
\crosscheck 相对半波损失符号取决于两束相减顺序；
若非正入射，应另计折射角及偏振，传播相位对应$2n_2e\cos r$，反射相位也须按菲涅耳系数确定。反射强弱只依模$2\pi$的相位差，不能遗漏$n_1\lambda_1=\lambda_0$。

\sourcetest{9.1.3}{202}{210}
\problemtext 图9.1.2：杨氏双缝光源S从两缝等距离位置向下移至S'，
观察中央明纹的位置变化及条纹间距。
\workings
S下移使下缝S$_2$相对S$_1$的入射光程缩短。
要恢复零光程差，屏上观察点需位于上方，使
S$_2$到该点的传播路径更长。
故中央明纹\textbf{上移}。
在旁轴近似下，两缝间距、波长、屏距均不变，
$\Delta x=\lambda D/d$仍不变，选\textbf{B}。
\crosscheck 这里只移动上游光源而没有改变双缝结构，
造成的是初相差的平移而非条纹密度变化。

\sourcetest{9.1.4}{202}{210}
\problemtext 白光照射双缝，分别用纯红和纯蓝滤光片遮住两缝，
判断屏上的干涉现象。
\workings
来自不同颜色的两缝光波频率不同；
长期平均的交叉干涉项
$\langle E_R E_B\rangle$随时间平均为零，
假设两滤光片的透射频谱互不重叠，且探测时间远大于拍频周期，因此没有两缝间稳定的红蓝相干条纹，故\textbf{不形成稳定双缝干涉条纹}，选\textbf{D}。
\crosscheck 可能见到各单缝的衍射照度分布，
不等于红蓝光之间会形成两套相干双缝干涉条纹。

\sourcetest{9.1.5}{202--203}{210--211}
\problemtext 图9.1.3：双缝原先在P点形成亮纹。
遮住S$_2$直接光后，将镜M放在其附近使反射光到P，
判断P点亮暗。
\workings
新光路与原S$_2$直达光路的几何长度一般不同，
而反射镜还引入与材料相关的反射相位。
若未进一步给镜的位置及光程差，不能确定
与S$_1$直达光在P处的相位是偶数还是奇数$\pi$。
故只能断言P点\textbf{不一定仍亮或变暗}，
选\textbf{C：无法确定}。
\crosscheck 即使理想金属镜近似产生固定反射相位，
光路的额外几何长度仍不可忽略。

\sourcetest{9.1.6}{203}{211}
\problemtext 空气劈尖的左端为棱边，上玻璃板绕该棱边
按原图约定逆时针微转，使劈角增大，判断反射等厚纹变化。
\workings
空气劈尖厚度$t(x)\simeq\alpha x$，
同级厚度$t_m$与角$\alpha$无关，
$x_m=t_m/\alpha$；邻纹距
\[
\Delta x=\frac{\lambda}{2\alpha}.
\]
$\alpha$增加使$\Delta x$变小，
同级纹向左侧棱边移动。
选\textbf{A（变密并向棱边移动）}。
\crosscheck 转角方向以源题的上板与左侧棱边图示为准；
若画图镜像，方向口述必须随图调整。

\sourcetest{9.1.7}{203}{211}
\problemtext 牛顿环透镜垂直向上缓慢移动而远离下面的平板玻璃，
问观察到的反射干涉环怎样变化。
\workings
空气膜厚度$t(r)=t_0+r^2/(2R)$，
固定级数暗环$2t(r)=m\lambda$，故
\[
r_m^2=R(m\lambda-2t_0).
\]
上移使$t_0$增大，固定级数的$r_m$减小，
环纹\textbf{向中心收缩}，选\textbf{B}。
\crosscheck 当$m\lambda<2t_0$时该级暗环消失，
不能继续取负的半径平方。

\sourcetest{9.1.8}{203}{211}
\problemtext 空气中独立放置折射率$n$透明薄膜，
波长$\lambda$的光垂直入射；反射光增强的最小正膜厚。
\workings
假设$n>1$（普通透明介质）且$\lambda$为真空波长；空气$\to$膜有一次半波相变，
膜$\to$空气反射无半波相变。
反射增强要求$2nt=(m+1/2)\lambda$，
取$m=0$，
\[
\boxed{t_{\min}=\frac{\lambda}{4n}},
\]
选\textbf{B}。
\crosscheck 此结果只针对\emph{空气--薄膜--空气}，
不能直接套用于膜下还有高折射率基底的减反膜问题。

\sourcetest{9.1.9}{203}{211}
\problemtext 玻璃基底$n_3=1.60$上镀MgF$_2$膜$n_2=1.38$，
空气$n_1=1.00$中500nm光垂直入射，求反射尽量弱时最小正膜厚。
\workings
由于$n_1<n_2<n_3$，两主要反射均有$\pi$相变，
相对反射相位损失抵消。
在无吸收、折射率固定的单层膜模型中，反射极小条件$2n_2t=\lambda/2$，所以
\[
\boxed{t_{\min}=\frac{500}{4(1.38)}\,\mathrm{nm}
\simeq90.6\,\mathrm{nm}},
\]
选\textbf{E}。
\crosscheck 相位相反并不足以保证反射完全抵消；完整多次反射模型在四分之一波膜厚处的反射率为$[(n_1n_3-n_2^2)/(n_1n_3+n_2^2)]^2$，仅当$n_2=\sqrt{n_1n_3}$才为零。
因此题面准确说法是“尽可能少反射”。

\sourcetest{9.1.10}{203}{211}
\problemtext 用劈尖干涉探查工件表面的局部凹凸，
原图9.1.5显示直条纹在缺陷区发生一个条纹间距的横向弯曲。
判断缺陷方向与高度/深度。
\workings
同一干涉级次要求空气膜厚度不变，
且相邻反射等厚纹的膜厚差为$\lambda/2$。
原图局部弯曲相当于一个条纹间距，
所以缺陷表面相对正常面有
\[
\boxed{|h|=\frac{\lambda}{2}}.
\]
根据原图劈尖由左到右增厚的斜率，
局部干涉纹向薄侧弯曲表示空气膜在缺陷处更厚，
因此工件表面应有\textbf{凹陷}，
选\textbf{C（凹陷深度$\lambda/2$）}。
\crosscheck 图中“条纹向哪侧弯”决定凹/凸符号。
如果原图扫描倒置或劈尖厚度方向被重画反转，
须重新逐线判定；相位级差$\lambda/2$不受此影响。

\sourcetest{9.1.11}{203--204}{211--212}
\problemtext 迈克耳孙干涉仪一条臂插入折射率$n$、
厚度$d$的薄片，求这一臂光程相对不插片时的改变。
\workings
在薄片正入射、替代同厚空气且不改变其余光路时，光线经过该薄片向镜面行进一次并返回一次，
每次较空气中的同厚光程增加$(n-1)d$，所以
\[
\boxed{\Delta\mathcal L=2(n-1)d},
\]
选\textbf{A}。
\crosscheck 迈克耳孙仪是往返光路，不能漏掉因子2。
% === END guide-selftest09-choices.tex ===

% === BEGIN guide-selftest09-blanks.tex blob c1f8ed80a507ed9d44a722a1cb0ba9a48be2cafa ===
% 自测题9：填空9.2.1—9.2.7。印刷204／PDF212页。
\subsection{自测题9：薄膜、劈尖与双缝填空}

\sourcetest{9.2.1}{204}{212}
\problemtext 单色光在折射率$n_1$介质中波长为$\lambda$，
垂直入射折射率$n>n_1$、厚度$e$薄膜，
膜两侧环境均为$n_1$；求上下表面反射光相遇的相位差。
\workings
真空波长$\lambda_0=n_1\lambda$。
上反射有$\pi$变相、下反射无反射半波损失；
膜内双程光程$2ne$对应相位$4\pi ne/\lambda_0$。
故
\[
\boxed{\Delta\phi=\frac{4\pi ne}{n_1\lambda}+\pi
\pmod{2\pi}}.
\]
\crosscheck 若定义两束相位差为上束减下束，
公式整体取负或把$\pi$换成$-\pi$，
但干涉可观测强度不变。

\sourcetest{9.2.2}{204}{212}
\problemtext 两相干点光源至A点等距离。
在S$_1$到A的光路插入厚度$e$、折射率$n$的玻璃片，
求A点相位差；若$\lambda=500$nm、$n=1.5$，
A恰为第四级亮纹中心，求$e$。
\workings
假设两源初相相同、光线正入射玻璃片，且玻璃片替代等厚空气，单程光程增加$(n-1)e$，
\[
\boxed{|\Delta\phi|=\frac{2\pi(n-1)e}{\lambda}}.
\]
第四级亮纹要求$(n-1)e=4\lambda$，
\[
\boxed{e=\frac{4(500\,\mathrm{nm})}{0.5}
=4000\,\mathrm{nm}=4.0\,\mu\mathrm m}.
\]
\crosscheck 本装置光线\emph{单程}经过玻璃，
不能误用迈克耳孙仪往返的因子2。

\sourcetest{9.2.3}{204}{212}
\problemtext 杨氏双缝在空气中条纹间距为1.00mm，
整个装置置于折射率$4/3$的水中，其他几何尺寸固定，
求新的条纹间距。
\workings
入水后频率不变，波长变为原来的$1/n$，
条纹间距$\Delta x=D\lambda/d$也缩短$1/n$：
\[
\boxed{\Delta x_{\rm water}=1.00/(4/3)\,\mathrm{mm}
=0.750\,\mathrm{mm}}.
\]
\crosscheck 题设固定几何距离；若实际透镜同时浸水并改变焦距，
需另计光学成像条件，但原题的杨氏屏距模型未涉及透镜。

\sourcetest{9.2.4}{204}{212}
\problemtext 空气中透明劈尖的楔角$\alpha=1.0\times10^{-4}$rad，
垂直入射波长700nm单色光，邻近同类反射干涉纹距0.25cm，
求膜折射率。
\workings
近轴小楔角$t(x)\simeq\alpha x$，相邻同类纹
厚度差$\lambda/(2n)$，
\[
\boxed{n=\frac{\lambda}{2\alpha\Delta x}
=\frac{700\times10^{-9}}{2(10^{-4})(2.5\times10^{-3})}
=1.40}.
\]
\crosscheck cm与mm换算：0.25cm=2.5mm，不能误作0.25mm。

\sourcetest{9.2.5}{204}{212}
\problemtext 牛顿环装置原为空气，某一第$k$级暗环半径$r_1$；
透镜与板之间填入折射率$n$的液体后
同一级暗环半径变为$r_2$，求$n$。
\workings
同一暗环级数和镜面曲率半径下
$r_k^2=k\lambda_0R/n$，
因此
\[
\boxed{n=\left(\frac{r_1}{r_2}\right)^2}.
\]
\crosscheck 要求$k>0$、接触处无额外间隙，且液体不改变反射相位关系；这里环级数应保持相同，
且默认填入液体不会改变反射相位的级数归属。

\sourcetest{9.2.6}{204}{212}
\problemtext 迈克耳孙仪反射镜平移0.620mm，
共有2300条纹移过参考点，求单色光波长（nm）。
\workings
镜面平移$\Delta d$使往返光程差改变$2\Delta d$，
每一$\lambda$改变对应一条纹迁移：
\[
\boxed{\lambda=\frac{2\Delta d}{N}
=\frac{2(0.620\times10^{-3})}{2300}
\simeq539.1\,\mathrm{nm}}.
\]
\crosscheck 漏掉因子2会使结果变成约270nm。

\sourcetest{9.2.7}{204}{212}
\problemtext 两束相干光各自光强都是$I_0$，
问干涉区内所能出现的最大光强。
\workings
假设两束光同频、完全相干且偏振方向相同，两束光场的振幅各为$\sqrt{I_0}$（忽略相同比例常数），
同相叠加时振幅翻倍，
\[
\boxed{I_{\max}=(\sqrt{I_0}+\sqrt{I_0})^2=4I_0}.
\]
\crosscheck 不能只把两个\emph{强度}相加为$2I_0$，
因为本题是相干振幅叠加。若偏振方向夹角为$\chi$，线偏振束的最大强度为$2I_0(1+|\cos\chi|)$，正交时无干涉项。
% === END guide-selftest09-blanks.tex ===

% === BEGIN guide-selftest09-solutions.tex blob a8595514cd8fbf133e445f2dc59fc01bccdd407e ===
% 自测题9：计算题9.3.1—9.3.5，证明题9.4.1；
% 源书印刷204—205，源PDF物理212—213页。
\subsection{自测题9：双缝、薄膜和牛顿环计算题}

\sourcetest{9.3.1}{204}{212}
\problemtext 杨氏双缝间距0.20mm、屏距1.0m，
入射400—760nm白光。
屏上离零级亮纹20mm的位置，哪些波长在该点最大增强？
\workings
几何光程差
\[
\Delta=\frac{dx}{D}
=\frac{(0.20\times10^{-3})(20\times10^{-3})}{1.0}
=4.0\times10^{-6}\,\mathrm m=4000\,\mathrm{nm}.
\]
第$m$级亮纹需$m\lambda=\Delta$，
于是$\lambda=4000/m$nm。
要求$400\le\lambda\le760$nm，
即$4000/760\le m\le4000/400$，
整数$m=6,7,8,9,10$。因此
\[
\boxed{\lambda\simeq666.7,\ 571.4,\ 500.0,\ 444.4,\ 400.0\ \mathrm{nm}}.
\]
\crosscheck 上述五个波长是教材旁轴近似的结果。若严格按给定几何（两缝为点源，屏上$x=20$ mm处），精确光程差为
\[
\Delta_{\rm ex}=\frac{2dx}{\sqrt{D^2+(x+d/2)^2}+\sqrt{D^2+(x-d/2)^2}}\approx3999.20022\,\mathrm{nm}.
\]
于是$m=10$对应$399.92002$ nm，略低于400 nm。若题给范围与几何均视为精确，合法级次仅$m=6,7,8,9$，波长约$666.533,571.314,499.900,444.356$ nm；若沿用旁轴公式与给定数据精度，则保留400 nm端点并注明近似。不能将两种精度下的计数混同。

\sourcetest{9.3.2}{205}{213}
\problemtext 两缝照射波长546.1nm的平面单色光，
屏距$D=2.00$m；
零级亮纹两侧的第5级亮纹相距12.0mm。
求(1)缝距；(2)从零级向一侧数至第20条亮纹的距离；
(3)若平面光斜入射，条纹间隔如何变化？
\workings
设相邻亮纹间距$\beta=\lambda D/d$。
两侧级次为$m=\pm5$，所以$10\beta=12.0$mm，
\[
\beta=1.20\,\mathrm{mm}.
\]
由此
\[
\boxed{d=\frac{\lambda D}{\beta}
=\frac{(546.1\times10^{-9})(2.00)}
{1.20\times10^{-3}}
\simeq9.102\times10^{-4}\,\mathrm m=0.910\,\mathrm{mm}}.
\]
若不计零级本身、所问为第20级亮纹，从零级起第20级的坐标大小为
\[
\boxed{x_{20}=20\beta=24.0\,\mathrm{mm}}.
\]
平面光斜入射添加固定的初始两缝相差$d\sin i$，
不改变相邻级次的几何路程差递增$\lambda$，
故在小角度理想条件下
\[
\boxed{\beta'=\beta=1.20\,\mathrm{mm}}.
\]
\crosscheck 若计零级为第一条，第20条为第19级、距离22.8 mm，原题计数措辞待核。屏上$\pm5$亮纹之间有\emph{10}个条纹间隔，
不是5个；斜入射移动整个图样的零点，而非必改变间距。

\sourcetest{9.3.3}{205}{213}
\problemtext 折射率$n_f=1.35$的透明膜镀于$n_s=1.50$玻璃上，
垂直入射光在600nm处反射相消、700nm处反射相长，
且600—700nm之间不存在其他反射极大或极小。
求薄膜厚度。
\workings
假设两波长之间膜折射率无显著色散，空气$\to$膜和膜$\to$玻璃两次反射都产生$\pi$相变，
相对半波损失为零。
记膜内双程光程$L=2n_ft$；
反射相长要求$L=m\lambda$，相消要求
$L=(k+\tfrac12)\lambda$。
600nm相消与700nm相长是\emph{相邻}极值，
故两者的等效级数相差半级：
\[
 L=(m+\tfrac12)(600\,\mathrm{nm})
 =m(700\,\mathrm{nm}).
\]
解$600m+300=700m$得$m=3$，
\[
\boxed{L=2100\,\mathrm{nm}},\qquad
\boxed{t=\frac{2100}{2(1.35)}\,\mathrm{nm}
\simeq777.8\,\mathrm{nm}}.
\]
\crosscheck 700nm对应相长第3级；
600nm对应相消的3.5级，
二者中间无其他半整数级极值。这里相消指反射极小，并非保证反射率严格为零。若误判一束反射
才有相变，答案可能错半个级次。

\sourcetest{9.3.4}{205}{213}
\problemtext 空气楔尖两板一端相接，单色光500nm垂直照射；
楔角$\alpha=2\times10^{-4}$rad。
向楔尖充入折射率$n=1.40$的液体，
求从棱边数起第5条\emph{明纹}移动距离和方向。
\workings
对玻璃板间低于玻璃折射率的介质膜，
两主要反射光有一次相对半波损失；
从接触处暗纹起第一明纹对应
$2nt=(1/2)\lambda$。
从棱边数起第5条明纹则有
$2nt=(5-\tfrac12)\lambda=4.5\lambda$。
位置为
$x_5(n)=4.5\lambda/(2n\alpha)$。
介质由空气变为$n=1.40$后
\[
\boxed{|\Delta x|
=\frac{4.5\lambda}{2\alpha}\left(1-\frac1n\right)
=\frac{4.5(500\,\mathrm{nm})}{4\times10^{-4}}
 \left(1-\frac1{1.4}\right)
\simeq1.607\,\mathrm{mm}}.
\]
移动方向为\textbf{向劈尖棱边（薄端）}。
\crosscheck 反射等厚纹接触处为暗，
若把第5条明纹当成“5个完整波长”将有半级误差。
液体只改变光学厚度，不改变题给楔角的几何条件。

\sourcetest{9.3.5}{205}{213}
\problemtext 原书图9.3.1中的牛顿环，平凸透镜曲率半径不变，
玻璃间原为空气$n=1.00$，
后来改为水$n'=1.33$。
求第$k$级暗环半径的相对缩小量$(r_k-r'_k)/r_k$。
\workings
在$k>0$、零接触间隙、反射相位关系不变的模型下，同一环级反射暗纹半径
$r_k^2=kR\lambda_0/n$，故
\[
 \frac{r'_k}{r_k}=\sqrt{\frac{n}{n'}}
 =\frac1{\sqrt{1.33}}.
\]
于是
\[
\boxed{\frac{r_k-r'_k}{r_k}
=1-\frac1{\sqrt{1.33}}\simeq0.1329=13.29\%}.
\]
\crosscheck 第$k$级数和曲率半径同时消去，
因此相对缩小率不依赖$k$或$R$。

\subsection{自测题9：干涉场的振幅证明}

\sourcetest{9.4.1}{205}{213}
\problemtext 杨氏双缝两光束在屏幕P点振动为
$E_1=E_0\sin\omega t$及
$E_2=E_0\sin(\omega t+\Phi)$，
证明合振幅可写成
$E_P=E_m\cos[(\pi d/\lambda)\sin\theta]$，
其中$E_m$是最大合振幅。
\workings
利用和差化积恒等式
\[
\begin{aligned}
E_P&=E_1+E_2\\
&=2E_0\cos\frac\Phi2\,
 \sin\left(\omega t+\frac\Phi2\right).
\end{aligned}
\]
因此有向振幅系数为$2E_0\cos(\Phi/2)$，
最大幅值$E_m=2E_0$。
在等偏振、远场夫琅禾费近似下，双缝缝距$d$，观测方向$\theta$时几何光程差
$\delta=d\sin\theta$，若双缝初相相同，
\[
\Phi=\frac{2\pi\delta}{\lambda}
=\frac{2\pi d}{\lambda}\sin\theta.
\]
于是
\[
\boxed{E_P=E_m\cos\left(\frac{\pi d}{\lambda}\sin\theta\right)}
\]
表示\emph{有向合振幅的系数}；
实际非负幅值应写
$|E_m\cos[(\pi d/\lambda)\sin\theta]|$。
\crosscheck 合成振动还包含时间相位
$\sin(\omega t+\Phi/2)$，
不能把有向振幅系数误认为完整瞬时电场。
% === END guide-selftest09-solutions.tex ===
